\section{Experiments \tagO}\label{sec:experiments}

\subsection{Setup}\label{sec:setup}

\textbf{Data and metrics.} RSTPReid~\cite{Zhu_2021} contains 20,505 images of 4,101 identities with two descriptions per image, split into 3,701 training, 200 validation and 200 test identities. The test protocol ranks a gallery of 1,000 images against 2,000 textual queries. We report Rank-$K$ accuracy (R@$K$, the fraction of queries whose first correct image appears in the top $K$), mean average precision (mAP) and mean inverse negative penalty (mINP)~\cite{jiang2023crossmodalimplicitrelationreasoning}. With 2,000 queries, R@1 is a binomial estimate whose standard error is about 1.1 points near 58\%; mAP averages over complete rankings and is less noisy.

\textbf{Training details.} We start from the publicly released CLIP ViT-B/16 weights~\cite{radford2021learningtransferablevisualmodels}. Training images are flipped, padded, randomly cropped and randomly erased, and captions are truncated to 77 tokens. We train with AdamW~\cite{loshchilov2019decoupledweightdecayregularization} (peak learning rate $10^{-4}$, weight decay 0.02), batch size \batchSize, gradient clipping at norm 5 and mixed precision. The learning rate rises linearly over 3 warm-up epochs and then follows a cosine decay to 1\% of its peak; one epoch takes about \sdmMinPerEp{} minutes on a single GPU. The SDM run uses a \sdmRounds-epoch schedule and the InfoNCE and MPCE runs a \nceRounds-epoch schedule, InfoNCE being stopped after \nceEpochs{} epochs. Warm-up is identical, and up to epoch 9 the learning rates stay close (at epoch 4, $\sdmLrFour$ vs.\ $\nceLrFour\times10^{-5}$; at epoch 9, $\sdmLrNine$ vs.\ $\nceLrNine\times10^{-5}$).

\textbf{Protocol.} We evaluate after every epoch. Like IRRA, we report the epoch with the best R@1 (\emph{peak}). Selecting an epoch on the test split is optimistic, so we also report the \emph{converged} value, the mean $\pm$ standard deviation over the last six epochs.

\begin{table*}[t]
\centering
\caption{Results on RSTPReid (text-to-image, \%). Rows marked IRRA are quoted from~\cite{jiang2023crossmodalimplicitrelationreasoning} (Tables 3 and 4 there; that ablation reports no mAP or mINP). Our rows are single runs at their peak epoch; converged values are in \cref{tab:ablation}. ``Stored'' is the size of the trained parameters in fp32.}
\label{tab:main}
\footnotesize
\setlength{\tabcolsep}{3.6pt}
\begin{tabular}{@{}llllrr ccccc@{}}
\toprule
& Method & Tuning & Loss & Trainable & Stored & R@1 & R@5 & R@10 & mAP & mINP\\
\midrule
\multirow{3}{*}{\rotatebox{90}{\scriptsize IRRA}}
& CLIP baseline & full & InfoNCE & 149.6\,M$^\ast$ & \fullMiB\,MiB & 54.05 & 80.70 & 88.00 & 43.41 & 22.31\\
& SDM only (Tab.~4, No.~2) & full & SDM & 149.6\,M$^\ast$ & \fullMiB\,MiB & 57.20 & 79.90 & 88.10 & -- & --\\
& IRRA (full model) & full & SDM+ID+IRR & $>$165\,M$^\dagger$ & $>$630\,MiB & \textbf{60.20} & \textbf{81.30} & 88.20 & \textbf{47.17} & \textbf{25.28}\\
\midrule
\multirow{3}{*}{\rotatebox{90}{\scriptsize Ours}}
& LoRA $r{=}4$ (ep.~\ncePeakEp) & LoRA & InfoNCE & 0.49\,M & \ckptMiB\,MiB & \ncePeakR & \ncePeakRfive & \ncePeakRten & \ncePeakM & \ncePeakMINP\\
& LoRA $r{=}4$ (ep.~\mpcePeakEp) & LoRA & MPCE & 0.49\,M & \ckptMiB\,MiB & \mpcePeakR & \mpcePeakRfive & \mpcePeakRten & \mpcePeakM & \mpcePeakMINP\\
& LoRA $r{=}4$ (ep.~\sdmPeakEp) & LoRA & SDM & 0.49\,M & \ckptMiB\,MiB & \sdmPeakR & \sdmPeakRfive & \textbf{\sdmPeakRten} & \sdmPeakM & \sdmPeakMINP\\
\bottomrule
\end{tabular}\\[2pt]
\parbox{\textwidth}{\scriptsize $^\ast$CLIP ViT-B/16 with position embeddings resampled to the $24{\times}8$ grid (\nBackbone{} parameters). $^\dagger$CLIP plus the 13.66\,M IRR encoder~\cite{jiang2023crossmodalimplicitrelationreasoning} and the \idHeadM\,M ID classifier; the MLM head is not counted.}
\end{table*}

\subsection{LoRA versus full fine-tuning (Q1)}\label{sec:main}

\Cref{tab:main} answers Q1. The fair comparison is IRRA's ablation row that, like ours, uses SDM alone. Full fine-tuning reaches 57.20\% R@1; LoRA reaches \sdmPeakR\% at its peak and \sdmConvR{}$\pm$\sdmConvRsd\% after convergence. The peak is \gapIrraSdmR{} points higher, about 1.5 standard errors, and the converged value matches, so with a single seed we read the two as on par rather than claim an improvement. LoRA gets there with \trainPct\% of the parameters and \ckptMiB\,MiB of trained weights instead of \fullMiB\,MiB. The complete IRRA model, which adds the ID loss and the IRR encoder, is ahead by \gapIrraR{} points of R@1 and \gapIrraM{} of mAP at the peak. Closing that gap would take extra objectives with their own parameters, the trade-off made by the PETL methods in \cref{sec:related}.

\subsection{Where does SDM's gain come from? (Q2)}\label{sec:ablation}

\begin{table}[t]
\centering
\caption{Loss ablation with LoRA $r{=}4$ (single seed, \%). Only the loss differs. Converged values average the last six epochs of each run's own schedule.}
\label{tab:ablation}
\footnotesize
\setlength{\tabcolsep}{3pt}
\begin{tabular}{@{}lllccc@{}}
\toprule
Loss & Target & Divergence & Peak R@1 (ep.) & Peak mAP & Conv.\ R@1\\
\midrule
InfoNCE & diagonal & CE & \ncePeakR{} (\ncePeakEp) & \ncePeakM & \nceConvR{}$\pm$\nceConvRsd\\
MPCE & identity & $\KL(q\Vert p)$ & \mpcePeakR{} (\mpcePeakEp) & \mpcePeakM & \mpceConvR{}$\pm$\mpceConvRsd\\
SDM & identity & $\KL(p\Vert q)$ & \textbf{\sdmPeakR} (\sdmPeakEp) & \textbf{\sdmPeakM} & \textbf{\sdmConvR}$\pm$\sdmConvRsd\\
\bottomrule
\end{tabular}
\end{table}

\begin{figure}[t]
\centering
\begin{tikzpicture}
\pgfplotsset{
  every axis/.append style={width=\columnwidth,height=3.55cm,xmin=0.5,xmax=30.5,
    tick label style={font=\scriptsize},label style={font=\scriptsize},
    grid=major,grid style={black!12},axis line style={black!45},tick style={black!45}},
  sdm/.style={cSDM,thick,mark=*,mark size=1.3pt},
  mpce/.style={cMPCE,thick,dashed,mark=triangle*,mark size=1.6pt,mark options={solid}},
  nce/.style={cNCE,thick,densely dotted,mark=square*,mark size=1.1pt,mark options={solid}}}
\begin{axis}[name=top,ylabel={R@1 (\%)},ymin=40,ymax=61,xticklabels={}]
\addplot[sdm] table[x=epoch,y=r1,col sep=comma]{generated/curve_sdm.csv};
\addplot[mpce] table[x=epoch,y=r1,col sep=comma]{generated/curve_mpce.csv};
\addplot[nce] table[x=epoch,y=r1,col sep=comma]{generated/curve_nce.csv};
\draw[cIRRA,thin,dashdotted] (axis cs:0.5,57.20) -- (axis cs:30.5,57.20) node[pos=0.98,above left,font=\tiny,inner sep=1pt]{IRRA full FT, SDM 57.20};
\draw[cIRRA,thin,dashdotted] (axis cs:0.5,54.05) -- (axis cs:30.5,54.05) node[pos=0.98,above left,font=\tiny,inner sep=1pt]{IRRA full FT, InfoNCE 54.05};
\end{axis}
\begin{axis}[at={(top.south)},anchor=north,yshift=-1.5mm,xlabel={epoch},ylabel={mAP (\%)},ymin=30,ymax=47,
  legend style={font=\scriptsize,at={(0.99,0.04)},anchor=south east,draw=black!30,fill=white,legend columns=3}]
\addplot[sdm] table[x=epoch,y=map,col sep=comma]{generated/curve_sdm.csv};
\addplot[mpce] table[x=epoch,y=map,col sep=comma]{generated/curve_mpce.csv};
\addplot[nce] table[x=epoch,y=map,col sep=comma]{generated/curve_nce.csv};
\legend{SDM,MPCE,InfoNCE}
\end{axis}
\end{tikzpicture}
\caption{Test R@1 and mAP after each epoch. MPCE and InfoNCE peak at the end of warm-up and then decline together; SDM keeps improving. Grey lines: IRRA's full fine-tuning results~\cite{jiang2023crossmodalimplicitrelationreasoning}.}
\label{fig:curves}
\end{figure}

\Cref{tab:ablation} and \cref{fig:curves} compare the three losses under identical LoRA training. SDM leads InfoNCE by \gapSdmNceR{} points of R@1 and \gapSdmNceM{} of mAP at the peak, close to the \gainIrraSdm{} points IRRA reports for full fine-tuning, so LoRA keeps SDM's benefit. The ablation also shows where that benefit comes from.

\textbf{Identity-aware targets contribute nothing measurable.} MPCE and InfoNCE share the same \nPairedEp-epoch schedule, so they can be compared epoch by epoch. Across these epochs MPCE's R@1 differs from InfoNCE's by \mpNceDiffR{}$\pm$\mpNceDiffRsd{} points on average (largest gap \mpNceDiffRmax) and its mAP by \mpNceDiffM{} (largest \mpNceDiffMmax), both well inside run-to-run noise. The ratio of their training losses stays between \lossRatioLo{} and \lossRatioHi, as \cref{sec:analysis} predicts: exactly 2 for one-hot rows, slightly more because of the \posFrac\% of rows with an extra positive, which confirms that the targets are almost always one-hot.

\textbf{The KL direction accounts for the whole gain.} SDM and MPCE share the target and differ only in the direction of the divergence, and SDM leads MPCE by \gapSdmMpceR{} points of R@1 and \gapSdmMpceM{} of mAP at the peak. The curves show how this gap arises. InfoNCE and MPCE peak at epoch 3, the end of warm-up, and then lose \nceDrop{} and \mpceDrop{} points of R@1 while their training losses keep falling (InfoNCE from \nceLossFirst{} to \nceLossLast), the usual sign of overfitting. SDM instead climbs until epoch \sdmPeakEp{} and ends within \sdmDrop{} points of its peak. Over epochs 4--9 the learning rates of the two schedules differ by at most \lrMaxRelDiff\%, yet the curves move in opposite directions (SDM from \sdmEpFourR{} to \sdmEpNineR, InfoNCE from \nceEpFourR{} to \nceEpNineR), so the schedule difference cannot explain the gap.

\textbf{Interpretation.} Under random sampling, IRRA's rationale for SDM, that it accounts for several positives per batch, applies to fewer than 2\% of rows and cannot explain a three-point gain. \Cref{eq:decomp,eq:grad} offer an explanation consistent with the curves: InfoNCE spends the largest gradients on the pairs it ranks worst, which on a small dataset with ambiguous descriptions leads it to memorise them, whereas SDM's bounded, MAE-like term down-weights those pairs and its entropy bonus discourages over-confident similarity distributions. We have not separated the two terms experimentally, so this remains a hypothesis, but the decomposition makes it testable: varying $\epsilon$ rescales $\lambda$ and with it the weight of the MAE term against the entropy bonus, and each term can be trained alone.

\subsection{Implications for federated training}\label{sec:fl}

\begin{table}[t]
\centering
\caption{What the trainable-part size implies for a federation of $K{=}\kClients$ clients (one per RSTPReid camera) over $R{=}\rRounds$ rounds. Arithmetic from the parameter counts (fp32), not a measurement; downloads from the server double the traffic.}
\label{tab:cost}
\footnotesize
\setlength{\tabcolsep}{4pt}
\begin{tabular}{@{}lrrr@{}}
\toprule
 & Full FT & LoRA $r{=}4$ & Ratio\\
\midrule
Trainable parameters & \nBackboneM\,M & 0.49\,M & \sizeRatio$\times$\\
Upload per client per round & \fullMiB\,MiB & \ckptMiB\,MiB & \sizeRatio$\times$\\
Upload per round ($K$ clients) & \upFullRoundGiB\,GiB & \upLoraRoundMiB\,MiB & \sizeRatio$\times$\\
Upload over $R$ rounds & \upFullTotalGiB\,GiB & \upLoraTotalGiB\,GiB & \sizeRatio$\times$\\
AdamW state per client & \optFullGiB\,GiB & \optLoraMiB\,MiB & \sizeRatio$\times$\\
\midrule
R@1, SDM only (peak / conv.) & 57.20 / -- & \sdmPeakR{} / \sdmConvR & \\
\bottomrule
\end{tabular}
\end{table}

\Cref{tab:cost} turns the parameter count into the quantities that dominate FL, assuming that every client takes part in every round and uploads its whole trainable part. With full fine-tuning each round would send \upFullRoundGiB\,GiB upstream, a heavy load for camera-side clients, whereas a whole LoRA round sends \roundVsOneFull$\times$ less than a single client's full-model upload. Each client's optimiser state shrinks by the same factor, which matters on edge hardware. Together with \cref{sec:main}, this is the premise a federated PEFT design for TBPS needs: the \sizeRatio$\times$ reduction costs no measurable accuracy against SDM-only full fine-tuning and \gapIrraR{} points of R@1 against the complete IRRA model.

The loss analysis carries over too. Unlike IRRA's ID classifier, whose rows refer to identities that other clients may never see, SDM has no parameters tied to the global identity set, so it can be optimised locally and averaged without modification. And since SDM's gain does not depend on several images of one person sharing a batch (\cref{sec:ablation}), there is no reason to expect it to vanish on clients with few images per identity. Whether these arguments hold under non-IID clients and the averaging of low-rank factors can only be settled by a federated experiment, which is beyond this report.
